\section{Reception: Doctor, Cardinal, Lion}

\subsection{From presbyter to Doctor of the Church}

Jerome's sanctity was received early and has never been questioned in the
Latin church: his feast stands on 30 September, and his authority as
interpreter of Scripture became proverbial within a century of his death. The
formal doctorate came in the Middle Ages: a decree of Boniface VIII (1298)
ordered the feasts of the four great Western teachers---Gregory the Great,
Ambrose, Augustine, and Jerome---``to be kept as doubles in the whole
Church,'' the act from which the title ``Doctor of the Church'' descends
(so the decretal \textit{Gloriosus} in the \textit{Liber Sextus}). The
Council of Trent, defining the canon and ordering a corrected edition,
made his composite Bible the Latin church's official text---a status
Benedict XVI's catechesis still summarizes---and the modern papacy has
returned to him three times at length: Benedict XV's encyclical
\textit{Spiritus Paraclitus} (1920), for the fifteenth centenary of his
death, proposing ``her `Greatest Doctor''' as the model of biblical study
and devotion; Benedict XVI's two general audiences of 7 and 14 November
2007, presenting the life (born c.~347, died 419/420) and drawing from it
the rule that reading Scripture is conversation with God; and Francis's
apostolic letter \textit{Scripturae Sacrae affectus} (30 September 2020),
for the sixteenth centenary, which retells the dream, celebrates the
``living and tender love'' for the written word of God, and reads Jerome's
iconography---the penitent of the desert and the scholar in the
study---as the two faces of one vocation. The Catechism's one direct
Jerome quotation (CCC 133) makes his sentence on scriptural ignorance the
Church's own exhortation.

\witnesslimit{Official reception is bounded by its objects. The doctorate
honors eminent doctrine and sanctity; the Tridentine decree concerned the
authenticity of a Bible, not the perfection of a translator; the papal
documents commend Jerome's love of Scripture and nowhere canonize his
polemical manners. None of these acts asserts the lion, the cardinalate,
the relic tradition, or either birth year as historical fact.}

\subsection{The lion}

The lion that pads through a thousand paintings entered the story late. No
ancient source knows it: not Jerome's own works, not his catalogue notice,
not Postumianus's eyewitness report, not the fifth-century chronicles (a
bounded negative result of this study's searches; see the appendix). The
full-blown tale appears in the medieval lives and classically in the
\work{Golden Legend} (thirteenth century): toward evening a lion ``came
halting suddenly in to the monastery''; the brethren fled, Jerome greeted it
as a guest, found ``the plant of the foot of the lion was sore hurt and
pricked with a thorn,'' healed it, and set the tame beast to guard the
monastery's donkey; merchants stole the donkey; the lion, suspected of
eating it, did the donkey's work in penance until it triumphantly herded
back the thieves' camel caravan (\work{Golden Legend}, Life of S. Jerome).
The same story-pattern---thorn, healing, grateful lion serving a
monastery---is told of the fifth-century abbot Gerasimus of the Jordan in
early Byzantine monastic literature, and the standard modern explanation,
argued in Eugene Rice's study of Jerome's iconography and repeated in
museum scholarship, is a transfer by confusion of names (Gerasimus /
Geronimus). The judgment of Jerome's own NPNF editor may stand for the
historical status: the desert descriptions ``no doubt gave rise to the
tradition that he was always attended by a lion\ldots\ With such traditions
a historical work must not be burdened'' (Fremantle, 1893). As reception,
however, the lion is eloquent: it gave art a single image for the tamed
ferocity that the documents show untamed.

\subsection{The cardinal's hat}

The \work{Golden Legend} also reports that at twenty-nine Jerome ``was
ordained cardinal priest in the church of Rome,'' and painters from the
fourteenth century onward robed him accordingly---the Thyssen catalogue's
\textit{Saint Jerome as a Cardinal} (1498) shows the full scarlet, with the
lion at his knee. The office is an anachronism: the cardinalate as a rank
did not exist in the fourth century, and the legend is a back-projection of
Jerome's real service as secretary and adviser to Pope Damasus, magnified
by centuries in which the pope's chief counselors were cardinals. The
red hat in the pictures is reception, not costume history; it testifies to
what later Rome felt it owed him, not to anything he wore.

\subsection{Relics and images}

Jerome was buried at Bethlehem; a later tradition, unverifiable now, holds
that his remains were translated to the basilica of Santa Maria Maggiore in
Rome, where an altar of St.~Jerome stands. The NPNF prolegomena reports it
in the proper mood: his remains ``are said to have been transferred\ldots\
to the Church of Santa Maria Maggiore at Rome, and miracles to have been
wrought at his tomb'' (Fremantle, 1893). No current authentication of relics
is asserted by this study, and none is required by the cult, which honors
the saint, not a verified skeleton. The iconography meanwhile split along
the line Francis's letter observes: D\"urer's scholar at his desk, skull and
lion in attendance; Leonardo's gaunt penitent with the stone; Caravaggio's
old man writing against time. Between them the painters preserved, better
than the legends, the two things the documents actually show: the desert
and the desk.

\subsection{The modern audit}

Modern scholarship has kept the achievement and cut away the varnish. The
standard biography (J.~N.~D. Kelly, 1975) and the current critical study
(Stefan Rebenich, 2002) present a Jerome who was a relentless
self-publicist as well as a genuine prodigy: the desert sojourn was
shorter, better supplied, and better connected than his rhetoric implies;
his Hebrew, though real and unprecedented, ``may well have been
exaggerated'' in his own telling and ``seems to have relied partly on the
hired assistance of Jewish scholars'' (so Rebenich, as quoted in the review
literature); and the Vulgate, as the previous section detailed, is a
composite in which his share, though decisive, is smaller than the legend
of the solitary translator. None of this scholarship doubts the scale of
the work; its effect has been to restore the collaborators---Paula,
Eustochium, Marcella, the Jewish teachers, the copyists---whom the lone
figure in the paintings had crowded out.
